\begin{claim}{D.39}{Proved}{A rank-two Haran involution admitting neither hyperbolic loop type}
Let $\epsilon,\kappa\in\widehat\Z$ have the coordinates
\[
 \epsilon_3=-1,\quad\epsilon_q=1\ (q\ne3),\qquad
 \kappa_2=1,\quad\kappa_q=0\ (q\ne2).
\]
There is an actual involution of the rank-two free
profinite group $H=\widehat F(a,b)$ given by
\[
 \alpha(a)=b^{-1}a^\epsilon b,\qquad
 \alpha(b)=b^{-1}a^\kappa.                                  \tag{D.39a}
\]
Its split extension has the explicit identification
\[
 H\rtimes_\alpha C_2\ \xrightarrow{\ \sim\ }\ N\amalg C_2,
 \qquad
 N=(\Z_3\rtimes\Z_2)\times\prod_{q\ne2,3}\Z_q,               \tag{D.39b}
\]
where parity in $\Z_2$ acts by inversion on $\Z_3$.
Its actual fixed subgroup $C_H(\alpha)$ is trivial.
The outer involution preserves no three-cusp
marking and no once-punctured-torus marking of $H$,
under any full profinite change of free basis.

For the natural once-punctured-torus marking
$\overline{\langle[a,b]\rangle}$, the peripheral
condition also fails: a quotient of order $21$
separates $\alpha([a,b])$ from every conjugate
of $[a,b]^{-\epsilon}$.
\end{claim}
\begin{proof}
The relations $\epsilon^2=1$ and
$\epsilon\kappa=\kappa$ give, by substitution
on the specified free generators,
\[
 \begin{split}
 \alpha^2(a)
  &=(a^{-\kappa}b)(b^{-1}ab)(b^{-1}a^\kappa)=a,\\
 \alpha^2(b)
  &=(a^{-\kappa}b)(b^{-1}a^{\epsilon\kappa}b)=b.
 \end{split}
\]
Thus the universal continuous endomorphism
defined by (D.39a) is an automorphism with
square one.
On full abelianization its matrix is
\[
       A=\begin{pmatrix}\epsilon&\kappa\\0&-1\end{pmatrix},
       \qquad \det A=-\epsilon,\qquad
       \operatorname{tr}A=\epsilon-1.                       \tag{D.39c}
\]
It is nonidentity on abelianization, so both
the automorphism and its outer class have
exact order two.

For the exact identification (D.39b), use
the normal form for $N$ verified in D.26:
\[
 N=\overline{\langle l,t\rangle},\qquad
 tlt^{-1}=l^\epsilon,\qquad t^2=l^\kappa.
\]
The map $l^u t^i\mapsto
(u_3,2u_2+i,(u_q)_{q\ne2,3})$, $i\in\{0,1\}$,
is its displayed isomorphism with the group
in (D.39b).  The parity character sends
$l$ to zero and $t$ to one.
In $N\amalg\langle r\mid r^2=1\rangle$
extend this character by $r\mapsto1$.
The substitutions
\[
             a\longmapsto l,\qquad
             b\longmapsto tr,\qquad
             s\longmapsto r
\]
define a homomorphism from
$H\rtimes_\alpha\langle s\rangle$ to that
free profinite product.  To check its
relations, note that $r=(tr)^{-1}t$,
so
\[
 rlr=(tr)^{-1}l^\epsilon(tr),\qquad
 r(tr)r=(tr)^{-1}l^\kappa.
\]
In the reverse direction send
$l\mapsto a$, $t\mapsto bs$, and $r\mapsto s$.
Here
$(bs)^2=b\alpha(b)=a^\kappa$ and
$(bs)a(bs)^{-1}=b\alpha(a)b^{-1}=a^\epsilon$.
The normal-form multiplication of $N$
therefore gives its continuous homomorphism
to the semidirect product, and the free
product property extends it over $r$.
The two composites fix every displayed
generator, proving the isomorphism.
In particular the parity kernel is exactly
the free profinite group on the stated
generators $a,b$, and the extension has
the actual section $\langle r\rangle$.

The free factor $\langle r\rangle$ is
malnormal by the full-profinite
free-factor intersection property
\cite[Theorem 5 of the inspected author
version]{A-HerfortZalesskii2008}.
Its centralizer in $N\amalg C_2$ is
therefore precisely $\langle r\rangle$.
Intersecting with the parity kernel
proves $C_H(\alpha)=1$.

In fact every involution $u\in H\rtimes_\alpha C_2$
has $C_H(u)=1$.  To prove this with the specified
index-two kernel, reduce (D.39c) modulo two.
The resulting matrix is
$\left(\begin{smallmatrix}1&1\\0&1\end{smallmatrix}\right)$.
Thus its dual action on $H^1(H,\F_2)$ is the
regular two-dimensional $\F_2[C_2]$ module $J_2$.
Its positive-degree $C_2$ cohomology vanishes:
both differentials of the periodic resolution are
$1+\alpha$, and their image equals their kernel,
both of dimension one.
Since $H$ has cohomological dimension one and
the extension is split, the Hochschild--Serre
spectral sequence gives
\[
       \dim H^n(H\rtimes_\alpha C_2,\F_2)=1
                        \qquad(n\geq3).                 \tag{D.39g}
\]
The full profinite Brown restriction theorem
\cite[Theorem 2 and \S1.1]{A-KochloukovaZalesskii2007},
applied exactly as in (D.37f), expresses this
cohomology as the product over involution classes.
Every class contributes
$1+\dim H^1(C_H(u),\F_2)$, because its normalizer
is $C_H(u)\times C_2$.
The existing section $\langle s\rangle$ therefore
accounts for the unique conjugacy class.
Conjugacy by the extension agrees with conjugacy
by $H$: for any involution $u$ the decomposition
$E=H\langle u\rangle$ absorbs its second factor,
which centralizes $u$.
Every involution is consequently $H$-conjugate
to $s$, and the already proved equality
$C_H(s)=1$ gives
\[
                  C_H(u)=1
           \quad\text{for every involution }u\in E.     \tag{D.39h}
\]

There is a basis-independent obstruction
to a three-cusp structure.  An involution
on three cusps has either the identity
permutation or a transposition.
In the former case its homology is a
scalar matrix $\lambda I$ and has trace
$2\lambda$, whose $2$-coordinate is
nonzero.  This contradicts the zero
$2$-coordinate of the trace in (D.39c).
For a transposition the reduced
permutation module has trace zero, and
multiplication by a common unit keeps
that trace zero.  This contradicts the
$3$-coordinate $-2$ of (D.39c).
Trace is unchanged by conjugating the
homology matrix under a change of free
basis.  Hence neither peripheral
permutation can occur.

There is also a basis-independent obstruction
to a once-punctured-torus marking.
Take any full profinite free basis $a',b'$ of
$H$ and put $c'=[a',b']$,
$L=\overline{\langle c'\rangle}$.
Preservation of the conjugacy class of $L$
would give the multiplier
$\nu=\det A=-\epsilon$.
The exact comparison is the one in
(D.33h): in every finite class-two
Heisenberg quotient the commutator
exponent is the determinant of the
abelianization matrix, so equality for
all moduli identifies these elements
of $\widehat\Z^\times$.
Changing the free basis conjugates the
matrix and leaves its determinant unchanged.

The normalizer $N_E(L)$ surjects onto
$E/H=C_2$, and its kernel is $L$ by A.7
transported through this free-basis isomorphism.
For a lift $v$ of the nonidentity element,
write
\[
  vc'v^{-1}=(c')^\nu,\qquad
  v^2=(c')^d,
  \qquad(\nu-1)d=0.                                    \tag{D.39i}
\]
Here $\nu_3=1$ and $\nu_q=-1$ for every
$q\ne3$.  Therefore $d_q=0$ for $q\ne3$,
including $q=2$, since multiplication by two
is injective on $\Z_q$.
Choose $j\in\widehat\Z$ with
$j_3=-d_3/2$ and $j_q=0$ otherwise.
The lift $u=(c')^jv$ then satisfies
\[
 u^2=(c')^{(1+\nu)j+d}=1,
 \qquad
 u(c')^{e_3w}u^{-1}=(c')^{e_3w}
                       \quad(w\in\Z_3),                 \tag{D.39j}
\]
where $e_3$ is the idempotent supported at
the prime three.
The lift $u$ has exact order two because
its image in $E/H$ is nonidentity.
The peripheral injection
$\widehat\Z\xrightarrow{\sim}L$ makes
$(c')^{e_3\Z_3}$ a nontrivial subgroup of
$C_H(u)$, contradicting (D.39h).
This excludes every once-punctured-torus
marking, including those obtained from
arbitrary full profinite free bases.

For an additional finite certificate in
the original basis, put
$c=[a,b]=aba^{-1}b^{-1}$.

Use the finite group
\[
 Q=\F_7\rtimes C_3,\qquad
 (v,i)(w,j)=(v+2^i w,\ i+j),
                                                               \tag{D.39d}
\]
where the second coordinate is modulo
three and $2$ has order three in
$\F_7^\times$.
The assignment
\[
                    a\mapsto(0,1),\qquad b\mapsto(1,1)
                                                               \tag{D.39e}
\]
is an epimorphism: $a,b$ have order
three, their commutator is $(1,0)$,
and these elements generate the
normal $\F_7$ and the $C_3$ quotient.
Since $\epsilon\equiv-1\pmod3$
and $\kappa\equiv0\pmod3$, (D.39a)
has the following exact images:
\[
 \begin{array}{c|ccccc}
 \text{word}&c&\alpha(a)&\alpha(b)&\alpha(c)&c^{-\epsilon}\\ \hline
 \text{image}&(1,0)&(5,2)&(3,2)&(1,0)&(6,0).
 \end{array}                                                   \tag{D.39f}
\]
The last entry uses
$(-\epsilon)_7=-1$.
Conjugacy on the normal $\F_7$ in
(D.39d) is multiplication by a
power of two.  The two relevant
orbits are exactly
\[
             \{(1,0),(2,0),(4,0)\},\qquad
             \{(3,0),(5,0),(6,0)\}.
\]
They are disjoint.  Thus the images
of $\alpha(c)$ and $c^{-\epsilon}$
are not conjugate, proving failure
of the natural torus cusp.

The program
\path{code/part_a/haran_torus_test.py}
executes this order-$21$ quotient
and an independent matrix check
inside $\operatorname{PSL}_2(\F_7)$.
For the latter it enumerates all
$168$ elements and uses
\[
 a=\begin{pmatrix}0&1\\6&1\end{pmatrix},
 \qquad b=\begin{pmatrix}0&2\\3&1\end{pmatrix}
 \quad\bmod\{\pm I\}.
\]
Both have order three,
$c=\left(\begin{smallmatrix}1&1\\3&4\end{smallmatrix}\right)$
has order seven, and
$\alpha(c)=
\left(\begin{smallmatrix}0&3\\2&2\end{smallmatrix}\right)$.
There are seven conjugators from
$c$ to this image and none from
$c^{-1}$.  The matrix calculation
agrees with (D.39f).
The finite calculation detects the original
cusp explicitly.  The argument
(D.39g)--(D.39j) supplies the exclusion
for every alternative marking.
The program also verifies the regular
mod-two module used in (D.39g).
\end{proof}
